% ============================================================================
% LANGENTWURF LE-013: E-Mail verstehen
% Profil: S (Senioren 65+) | Tier: 2 (Standard)
% ============================================================================

\documentclass[11pt,a4paper]{article}
\usepackage[utf8]{inputenc}
\usepackage[T1]{fontenc}
\usepackage[ngerman]{babel}
\usepackage{geometry}
\usepackage{fancyhdr}
\usepackage{lastpage}
\usepackage{xcolor}
\usepackage{tcolorbox}
\usepackage{enumitem}
\usepackage{longtable}
\usepackage{hyperref}
\usepackage{amssymb}

\geometry{left=2.5cm, right=2.5cm, top=2.5cm, bottom=2.5cm}
\definecolor{fach}{HTML}{b35c00}
\definecolor{gray}{HTML}{666666}
\definecolor{successgreen}{HTML}{2a7a4b}
\definecolor{warnred}{HTML}{c62828}

\tcbuselibrary{skins,breakable}
\newtcolorbox{scorebox}{ colback=white, colframe=fach, fonttitle=\bfseries, title=Didaktik-Score, breakable }

\pagestyle{fancy}
\fancyhf{}
\fancyhead[L]{\small\textcolor{gray}{Digitale Grundlagen -- 013 E-Mail verstehen}}
\fancyhead[R]{\small\textcolor{gray}{LE Tier 2 / Profil S}}
\fancyfoot[C]{\small\thepage{} / \pageref{LastPage}}
\renewcommand{\headrulewidth}{0.4pt}

\begin{document}

\begin{titlepage}
    \centering
    \vspace*{2cm}
    {\Large\textcolor{gray}{Digitale Grundlagen -- Senioren 65+}}\\[2cm]
    {\Huge\textcolor{fach}{\textbf{Langentwurf}}}\\[0.5cm]
    {\large Tier 2 | Profil S}\\[2cm]
    {\LARGE\textbf{013 -- E-Mail verstehen}}\\[0.5cm]
    {\large Modul C: Internet \& Kommunikation}\\[2cm]
    \begin{tabular}{rl}
        \textbf{Zeitrahmen:} & 60--75 Minuten \\
        \textbf{Vorkenntnisse:} & Safari, Suchmaschinen \\
    \end{tabular}
    \vfill
    {\normalsize Stand: \today}
\end{titlepage}

\tableofcontents
\newpage

\section{Bedingungsanalyse}

\subsection{Ausgangssituation}
\begin{itemize}
    \item E-Mail ist oft eingerichtet, aber kaum verstanden
    \item Angst vor ``falschen Klicks''
    \item Posteingang ist überfüllt (nie gelöscht)
    \item Anhänge werden nicht gefunden oder geöffnet
    \item Spam und wichtige Mails nicht unterschieden
\end{itemize}

\subsection{Typische Probleme}
\begin{itemize}
    \item ``Ich finde die E-Mail nicht mehr''
    \item ``Da war ein Anhang, aber wo?''
    \item ``Meine Mails sind alle weg'' (anderer Ordner)
    \item ``Ist diese Mail echt oder Betrug?''
\end{itemize}

\section{Sachanalyse}

\subsection{Kernkonzepte}
\begin{enumerate}
    \item \textbf{E-Mail-Adresse:} Wie eine Postadresse im Internet (name@anbieter.de)
    \item \textbf{Posteingang:} Hier landen neue Nachrichten
    \item \textbf{Gesendet:} Kopien deiner verschickten Mails
    \item \textbf{Papierkorb:} Gelöschte Mails (noch rettbar)
    \item \textbf{Spam/Werbung:} Automatisch aussortierte verdächtige Mails
    \item \textbf{Anhang:} Dateien, die an die Mail angehängt sind (Fotos, Dokumente)
\end{enumerate}

\subsection{Alltagsanalogien}
\begin{tabular}{|l|p{8cm}|}
\hline
\textbf{Begriff} & \textbf{Analogie} \\
\hline
E-Mail-Adresse & Deine Postadresse -- eindeutig, damit Briefe ankommen \\
Posteingang & Briefkasten -- hier landet die Post \\
Gesendet & Durchschläge der Briefe, die du geschickt hast \\
Spam & Die Werbezettel, die du normalerweise wegwirfst \\
Anhang & Etwas, das du dem Brief beilegst (Foto, Dokument) \\
Papierkorb & Der Mülleimer -- Inhalt ist noch drin, bis er geleert wird \\
\hline
\end{tabular}

\section{Didaktische Analyse}

\subsection{Stellung in der Reihe}
Grundlegende Kommunikationseinheit. E-Mail wird für viele Dienste benötigt (Registrierung, Bestätigungen).

\subsection{Didaktische Reduktion}
\textbf{Ausgeklammert:} E-Mail-Einrichtung, Signaturen, Regeln/Filter, mehrere Konten.

\textbf{Fokus:} Die Mail-App verstehen, Mails finden und lesen, Anhänge öffnen.

\section{Lernziele}

\begin{tabular}{|c|p{7cm}|c|c|}
\hline
\textbf{Nr.} & \textbf{Die Teilnehmer können...} & \textbf{Stufe} & \textbf{Niveau} \\
\hline
FZ1 & die Mail-App öffnen und den Posteingang finden & Anwenden & Basis \\
FZ2 & eine E-Mail öffnen und lesen & Anwenden & Basis \\
FZ3 & die verschiedenen Ordner erklären & Verstehen & Basis \\
FZ4 & einen Anhang öffnen und speichern & Anwenden & Basis \\
FZ5 & eine verdächtige Mail erkennen (Spam, Phishing) & Anwenden & Vertiefung \\
\hline
\end{tabular}

\section{Verlaufsplanung}

\begin{longtable}{|p{1cm}|p{2cm}|p{5cm}|p{3cm}|p{2cm}|}
\hline
\textbf{Zeit} & \textbf{Phase} & \textbf{Inhalt} & \textbf{TN-Aktivität} & \textbf{Material} \\
\hline
0--10 & Einstieg & ``Wer nutzt E-Mail? Wofür?'' Brief vs. E-Mail & Austauschen & -- \\
\hline
10--25 & App kennen & Mail-App öffnen, Bereiche erklären (Postfächer, Mails) & Mitmachen & S.1 \\
\hline
25--40 & Ordner & Posteingang, Gesendet, Papierkorb, Spam erklären & Eigene Ordner anschauen & S.1 \\
\hline
40--55 & Anhänge & Was ist ein Anhang? Öffnen, Speichern (Fotos, PDFs) & Üben & S.2 \\
\hline
55--70 & Sicherheit & Spam erkennen, verdächtige Mails, Wiederholung 010 & Beispiele bewerten & S.2-3 \\
\hline
\end{longtable}

\section{Lösungen}

\textbf{Mail-App auf iPhone:}
\begin{itemize}
    \item Öffnen: Mail-Symbol (weißer Umschlag auf blauem Grund)
    \item Links: Postfächer-Übersicht
    \item Mitte: Liste der Mails im gewählten Ordner
    \item Antippen: Mail wird rechts/unten angezeigt
\end{itemize}

\textbf{Anhang öffnen:}
\begin{itemize}
    \item Anhänge erscheinen unter dem Mailtext
    \item Antippen öffnet eine Vorschau
    \item Teilen-Symbol zum Speichern oder Weiterleiten
\end{itemize}

\textbf{Wichtige Ordner:}
\begin{itemize}
    \item \textbf{Eingang:} Neue und gelesene Mails
    \item \textbf{Gesendet:} Deine verschickten Mails
    \item \textbf{Entwürfe:} Angefangene, noch nicht gesendete Mails
    \item \textbf{Papierkorb:} Gelöschte Mails (30 Tage aufbewahrt)
    \item \textbf{Spam/Werbung:} Automatisch aussortiert -- manchmal prüfen!
\end{itemize}

\section{Didaktik-Score}

\begin{scorebox}
\begin{tabular}{|c|l|c|c|p{3.5cm}|}
\hline
\# & Dimension & Gew. & Score & Notiz \\
\hline
1 & Differenzierung & 15\% & 8/10 & Basis stark, Vertiefung ok \\
2 & Taxonomie & 10\% & 9/10 & Verstehen + Anwenden \\
3 & Alltagsrelevanz & 10\% & 10/10 & E-Mail ist unverzichtbar \\
4 & Aktivierung & 10\% & 9/10 & Eigene Mails anschauen \\
5 & Scaffolding & 10\% & 9/10 & Gute Briefkasten-Analogie \\
6 & Feedback & 10\% & 8/10 & Checkliste vorhanden \\
7 & Sprachliche Klarheit & 10\% & 9/10 & Einfache Erklärungen \\
8 & Motivation & 10\% & 9/10 & Praktischer Nutzen klar \\
9 & Barrierefreiheit & 10\% & 9/10 & Große Darstellung \\
10 & Praktikabilität & 5\% & 9/10 & 70 Min gut \\
\hline
\multicolumn{3}{|r|}{\textbf{GESAMT:}} & \textbf{8.95/10} & Knapp unter Freigabe \\
\hline
\end{tabular}
\end{scorebox}

\end{document}
